\chapter{Proof Details for the Impossibility
Theorem}

This appendix records the formal development behind Chapter 9: the model
and the algorithm class (§B.1), the master trade-off inequality with its
proof (§B.2), the rare-resource cell construction and its closed-form
constants (§B.3), the exact affine law converting advice error into
matching ratio (§B.4), the reduction to tolerant distribution testing
and the sample-complexity barrier (§B.5), the assembled theorem (§B.6),
and a plain statement of what remains open (§B.7). The numerical
verifications quoted here are reproduced by the snippets referenced in
Appendix A.6.

\section{Setup and the algorithm
class}

\textbf{Instances.} Known-i.i.d. online bipartite matching (Chapter 3):
a type graph over \(r\) online types and a set of offline resources;
\(n\) arrivals drawn i.i.d. from a type distribution \(p\) over \([r]\);
\(\mathrm{OPT}(I)\) is a maximum matching of the realized instance
\(I\).

\textbf{Advice.} A type histogram \(q\in\Delta([r])\) (the Choo/BEM
prediction object). \emph{Following} the advice --- the \emph{mimic}
policy \(\mathrm{Mimic}(q)\) --- builds a maximum matching on the advice
graph and routes each arrival to its type's advice partner (Chapter 3).

\textbf{Baseline.} The advice-free Ranking algorithm \(B\), with
\(\mathbb E[B]/\mathbb E[\mathrm{OPT}] = \rho_{\mathrm{base}}\) on the
family under study.

\textbf{The test-and-fallback class \(\mathcal A_k\).} An algorithm
\(A\in\mathcal A_k\) observes the advice \(q\) and the prefix
\(X_{1:k}\) (the types of the first \(k\) arrivals), outputs a ---
possibly randomized --- decision \(D\in\{\mathrm F,\mathrm R\}\)
\emph{by any measurable rule}, and then plays \(\mathrm{Mimic}(q)\) on
arrivals \(k+1,\dots,n\) if \(D=\mathrm F\) and \(B\) otherwise. The
class covers every decision rule, not only the empirical-\(\ell_1\)
threshold used by the concrete algorithms.

\textbf{Consistency and robustness.} For an instance distribution
\(\mathcal D\) with advice \(q\), write
\(R(\mathcal D)=\mathbb E[\mathrm{ALG}]/\mathbb E[\mathrm{OPT}]\).
Consistency is \(R\) under good advice; robustness is \(R\) under bad
advice (Chapter 3).

\textbf{Sublinearity (A0).} \(k=o(n)\). The prefix's own contribution to
the matching is at most \(k=o(n)\), shifting every ratio by
\(O(k/n)=o(1)\); we absorb this into the \(o(1)\) terms.

\section{The master trade-off
inequality}

\begin{quote}
\textbf{Lemma B.1 (master trade-off).} Let \(G\), \(\mathrm{Bd}\) be two
instance distributions sharing the \emph{same} advice \(q\), with prefix
laws \(\mathcal L_G,\mathcal L_{\mathrm{Bd}}\) on \(X_{1:k}\), and let
\(\gamma_k := \mathrm{TV}(\mathcal L_G,\mathcal L_{\mathrm{Bd}})\).
Suppose following \(q\) gains \(\delta\) under \(G\) and loses
\(\Delta\) under \(\mathrm{Bd}\) relative to the baseline:
\[\mathbb E_G[\mathrm{Mimic}]/\mathrm{OPT} \ge \rho_{\mathrm{base}}+\delta,
\qquad \mathbb E_{\mathrm{Bd}}[\mathrm{Mimic}]/\mathrm{OPT} \le \rho_{\mathrm{base}}-\Delta,\]
while \(B\) attains \(\rho_{\mathrm{base}}\pm o(1)\) under both.
Parameterize an \(A\in\mathcal A_k\) by \(\eta_c\), the fraction of the
upside it forgoes under \(G\) (captured upside \(=(1-\eta_c)\,\delta\)),
and \(\eta_r\), its robustness loss as a fraction of \(\Delta\) under
\(\mathrm{Bd}\) (loss \(=\eta_r\,\Delta\)). Then
\[(1-\eta_c)\;\le\;\eta_r+\gamma_k+o(1). \tag{$\star$}\]
\end{quote}

\emph{Proof.} Condition on the decision \(D\) under \(G\):
\[\mathbb E_G[\mathrm{ALG}]/\mathrm{OPT}
= P_G(\mathrm F)\cdot\frac{\mathbb E_G[\mathrm{Mimic}]}{\mathrm{OPT}}
+ P_G(\mathrm R)\cdot\rho_{\mathrm{base}} \pm o(1)
\;\ge\; \rho_{\mathrm{base}}+\delta\,P_G(\mathrm F)-o(1).\] The captured
upside is \((1-\eta_c)\delta\) by definition, so
\(P_G(\mathrm F)\ge 1-\eta_c-o(1)\). Symmetrically under
\(\mathrm{Bd}\), \(\mathbb E_{\mathrm{Bd}}[\mathrm{ALG}]/\mathrm{OPT}
= \rho_{\mathrm{base}}-\Delta\,P_{\mathrm{Bd}}(\mathrm F)\pm o(1)\), so
the robustness loss is \(\Delta\,P_{\mathrm{Bd}}(\mathrm F)\pm o(1)\),
giving \(P_{\mathrm{Bd}}(\mathrm F)=\eta_r+o(1)\). Finally, \(D\) is a
(randomized) function of \((X_{1:k},q)\) and \(q\) is identical across
\(G\) and \(\mathrm{Bd}\), so by the coupling characterization of total
variation, \(|P_G(\mathrm F)-P_{\mathrm{Bd}}(\mathrm F)|\le \gamma_k\).
Chaining the three displays yields
\(1-\eta_c-o(1)\le P_G(\mathrm F)\le P_{\mathrm{Bd}}(\mathrm F)+\gamma_k
= \eta_r+\gamma_k+o(1)\), which is \((\star)\). \(\blacksquare\)

When the prefix is uninformative (\(\gamma_k\to0\)), \((\star)\) forbids
\(\eta_c\to0\) and \(\eta_r\to0\) simultaneously: the algorithm cannot
be both consistent and robust. The remaining work is to construct a
family on which \(\gamma_k=o(1)\) \emph{coexists} with meaningful stakes
\(\delta,\Delta=\Theta(1-\rho_{\mathrm{base}})\).

\textbf{Remark (why a two-point argument is not enough).} For a single
pair \((G,\mathrm{Bd})\) whose per-arrival laws differ, the stakes
themselves are capped by the distinguishability:
\(\delta,\Delta = O(\gamma_k)\), so \((\star)\) alone gives only a
vacuous bound. To obtain \(\delta,\Delta=\Theta(1)\) \emph{while}
\(\gamma_k=o(1)\), the difference between \(G\) and \(\mathrm{Bd}\) must
be spread across \(\Theta(r)\) types, each perturbed below the per-type
sampling resolution of a length-\(k\) prefix --- which is precisely the
structure of the \emph{tolerant-testing} hard instances invoked in §B.5.
This is also why a low-dimensional construction fails (§B.7).

\section{The rare-resource cell}

The construction is assembled from independent \emph{cells}, each
computable in closed form.

\textbf{The cell.} Two offline resources \(\{a,b\}\). Arrivals, in
order: a \emph{flexible} request \(F\) with neighborhood \(\{a,b\}\),
which always arrives and must be routed \emph{before} the specialist is
seen; then at most one \emph{specialist} --- type \(\alpha\)
(neighborhood \(\{a\}\)) with probability \(\theta s\), type \(\beta\)
(neighborhood \(\{b\}\)) with probability \(\theta(1-s)\), or no
specialist with probability \(1-\theta\). Here \(\theta\in(0,1]\) is the
\textbf{contention rate} and \(s\in[0,1]\) the \textbf{bias}. The advice
specifies a predicted bias \(\hat s\); following it routes \(F\) to
protect the predicted-majority specialist (\(\hat s>\tfrac12\): route
\(F\to b\), saving \(a\) for the likelier \(\alpha\); symmetrically
otherwise).

\begin{quote}
\textbf{Lemma B.2 (single-cell constants).} Per cell, in expectation:

{ % do not increment counter
\begin{longtable}[]{@{}
  >{\raggedright\arraybackslash}p{(\linewidth - 2\tabcolsep) * \real{0.5000}}
  >{\raggedright\arraybackslash}p{(\linewidth - 2\tabcolsep) * \real{0.5000}}@{}}
\toprule\noalign{}
\begin{minipage}[b]{\linewidth}\raggedright
quantity
\end{minipage} & \begin{minipage}[b]{\linewidth}\raggedright
value
\end{minipage} \\
\midrule\noalign{}
\endhead
\bottomrule\noalign{}
\endlastfoot
\(\mathrm{OPT}\) & \(1+\theta\) \\
Baseline (route \(F\) uniformly) & \(1+\theta/2\) \\
Mimic, advice direction \emph{right}
(\(\operatorname{sign}(\hat s-\tfrac12)=\operatorname{sign}(s-\tfrac12)\))
& \(1+\theta\max(s,1-s)\) \\
Mimic, advice direction \emph{wrong} & \(1+\theta\min(s,1-s)\) \\
\end{longtable}
}

Hence the per-cell advantage of following over the baseline is
\(\pm\,\theta\,|s-\tfrac12|\) (sign \(=\) is the advice direction
correct), and the \(\ell_1\) distance between the cell's specialist
distribution \(p=(\theta s,\ \theta(1-s),\ 1-\theta)\) and the advice
\(q=(\theta\hat s,\ \theta(1-\hat s),\ 1-\theta)\) is
\(\ell_1(p,q)=2\theta\,|s-\hat s|\).
\end{quote}

\emph{Proof.} Routing \(F\to a\) matches \(F\) always and the specialist
unless it is \(\alpha\) (whose only neighbor \(a\) is taken):
\(\theta s\cdot 1+\theta(1-s)\cdot 2+(1-\theta)\cdot 1 = 1+\theta(1-s)\).
Symmetrically \(F\to b\) yields \(1+\theta s\). The baseline averages
the two, \(1+\theta/2\). The better route is \(F\to b\) iff
\(s>\tfrac12\), giving \(1+\theta\max(s,1-s)\) for the right direction
and \(1+\theta\min(s,1-s)\) for the wrong one. \(\mathrm{OPT}=2\) when a
specialist arrives (probability \(\theta\)), else \(1\), totalling
\(1+\theta\). The \(\ell_1\) computation is a three-coordinate sum whose
\((1-\theta)\) coordinate cancels. \(\blacksquare\)

Two quantities separate cleanly, which is the point of the design: the
\emph{magnitude} of the advantage, \(\theta|s-\tfrac12|\), depends only
on the truth's signal (how contested the cell is), while its \emph{sign}
depends only on whether the advice's direction is right. Deciding
whether the advice is net-helpful is therefore exactly the problem of
recovering per-cell \emph{directions} --- a Bernoulli discrimination
needing \(\approx 1/|s-\tfrac12|^2\) samples \emph{of that cell}.
(Numerical check, \(4\times10^5\) trials per cell: at \(\theta=0.6\),
\(s=0.7\) the simulated advantage is \(\pm0.119\) against the formula's
\(\pm0.12\); all five identities of Lemma B.2 match to three decimals
--- Appendix A.6.)

\section{The aggregate construction and the exact affine
law}

\textbf{The family.} \(m=\Theta(n)\) independent cells, each with a
\emph{distinct} specialist type-pair \((\alpha_i,\beta_i)\), so the type
support is \(r=\Theta(n)\) and each specialist type is seen \(O(1)\)
times among \(n\) arrivals. Cell \(i\) has contention \(\theta_i\) and
truth bias \(s_i\) with signal \(\varepsilon_i=|s_i-\tfrac12|\); the
advice predicts, per cell, a direction with the same magnitude
(\(\hat s_i=\tfrac12\pm\varepsilon_i\), correct or wrong side). Because
cells are disjoint, per-cell quantities aggregate additively.

\begin{quote}
\textbf{Lemma B.3 (exact affine conversion law).} Let
\(\rho_{\mathrm{perfect}}=\bigl(\sum_i(1+\theta_i\max(s_i,1-s_i))\bigr)/\mathrm{OPT}\)
be the all-directions-correct (consistency ceiling) ratio, with
\(\mathrm{OPT}=\sum_i(1+\theta_i)\). Then, for any advice with per-cell
magnitudes matching the truth,
\[\mathbb E[\text{follow-ratio}] \;=\; \rho_{\mathrm{perfect}}-\tfrac12\,\ell_1(p,q),\]
and hence the break-even advice error is
\(\ell_1^* = 2(\rho_{\mathrm{perfect}}-\rho_{\mathrm{base}})\), where
\(\rho_{\mathrm{base}}=\bigl(\sum_i(1+\theta_i/2)\bigr)/\mathrm{OPT}\).
\end{quote}

\emph{Proof.} By Lemma B.2, cell \(i\) contributes
\(1+\theta_i\max(s_i,1-s_i)\) to the followed matching if its advice
direction is correct and \(1+\theta_i\min(s_i,1-s_i)\) if wrong; the
difference is \(2\theta_i\varepsilon_i\). Letting \(W\) be the set of
wrong-direction cells, \[\text{followed matching}
= \sum_i\bigl(1+\theta_i\max(s_i,1-s_i)\bigr)-\sum_{i\in W}2\theta_i\varepsilon_i.\]
For the \(\ell_1\): the flexible and no-specialist coordinates cancel,
and cell \(i\) contributes \(2\theta_i|s_i-\hat s_i|\), which is \(0\)
if correct and (opposite side, same magnitude, so
\(|s_i-\hat s_i|=2\varepsilon_i\)) \(4\theta_i\varepsilon_i\) if wrong.
The type normalizer is \(N=\sum_i(1+\theta_i)=\mathrm{OPT}\), so
\(\ell_1(p,q)=\bigl(\sum_{i\in W}4\theta_i\varepsilon_i\bigr)/\mathrm{OPT}
=2\bigl(\sum_{i\in W}2\theta_i\varepsilon_i\bigr)/\mathrm{OPT}\).
Substituting, \[\mathbb E[\text{follow-ratio}]
=\frac{\text{followed matching}}{\mathrm{OPT}}
=\rho_{\mathrm{perfect}}-\tfrac12\,\ell_1(p,q). \qquad\blacksquare\]

The followed matching is a sum of \(m\) independent bounded cell
contributions, so it concentrates around its mean at rate
\(O(1/\sqrt m)\); the affine law holds with high probability, not only
in expectation. (Numerical verification at \(\theta=0.6\),
\(\varepsilon=0.3\), \(m=4000\) --- \(\rho_{\mathrm{perfect}}=0.924\),
\(\rho_{\mathrm{base}}=0.812\) --- matches the law to three decimals at
every wrong-fraction \(\varphi\in\{0,0.2,0.5,0.8,1\}\), with the
simulated break-even \(0.223\) against the predicted
\(\ell_1^*=2(0.924-0.812)=0.224\); Appendix A.6.)

\textbf{Consequence: the stakes and the gap are clean, and the good side
is a ball.} Define ``following gains \(\ge\delta\)''
\(\iff \ell_1\le \ell_1^*-2\delta =: a\) and ``following loses
\(\ge\Delta\)'' \(\iff \ell_1\ge \ell_1^*+2\Delta =: b\). Choosing,
e.g., \(\delta=\Delta=(\rho_{\mathrm{perfect}}-\rho_{\mathrm{base}})/4\)
gives \(a=\ell_1^*/2>0\), \(b=3\ell_1^*/2\), and gap \(b-a=\ell_1^*\)
--- all \(\Theta(1)\) constants independent of \(n\) (set by the fixed
per-cell \(\theta,\varepsilon\); only the number of cells grows).
Crucially \(a>0\): good advice is a \(\Theta(1)\)-radius \emph{ball}
around the truth, not the single point \(p=q\). This is what places the
decision problem in the \emph{tolerant} testing regime (§B.5).

\section{Reduction to tolerant identity
testing}

\begin{quote}
\textbf{Lemma B.4 (algorithm \(\Rightarrow\) tester).} On the family of
§B.4, fix the advice \(q\). Suppose \(A\in\mathcal A_k\) is
\((1-\eta_c)\)-consistent whenever \(\ell_1(p,q)\le a\) and loses at
most \(\eta_r\Delta\) whenever \(\ell_1(p,q)\ge b\). Then the decision
\(D\) of \(A\), viewed as a function of the \(k\) i.i.d. samples
\(X_{1:k}\sim p\), distinguishes \(\{\ell_1(p,q)\le a\}\) from
\(\{\ell_1(p,q)\ge b\}\) with error at most \(\eta_c+\eta_r+o(1)\).
\end{quote}

\emph{Proof sketch.} Consistency forces \(D=\mathrm F\) with high
probability when \(\ell_1\le a\) (otherwise the algorithm forgoes the
\(\delta\) gain, contradicting \((1-\eta_c)\)-consistency); robustness
forces \(D=\mathrm R\) with high probability when \(\ell_1\ge b\)
(otherwise it eats the \(\Delta\) loss). So \(D\) itself is the tester,
and its error probabilities are bounded by the consistency and
robustness slacks via the same conditioning as in Lemma B.1.
\(\blacksquare\)

The key point is that the prefix \(X_{1:k}\) is \emph{literally} \(k\)
i.i.d. samples from the type distribution \(p\) (the known-i.i.d.
model), and \(D\) is an arbitrary function of them --- so a
sample-complexity lower bound for the testing problem rules out
\emph{every} rule in \(\mathcal A_k\), not only the empirical-\(\ell_1\)
threshold. This is the step that makes the result algorithm-independent
and distinguishes it from the constructive baseline-coupling inside Choo
et al.'s algorithm (Chapter 9).

\begin{quote}
\textbf{Theorem B.5 (tolerant identity testing; cited).} Distinguishing
\(\ell_1(p,q)\le\varepsilon_1\) from \(\ell_1(p,q)\ge\varepsilon_2\) for
a known \(q\) over support size \(r\), with error \(\le 1/3\), requires
\[k \;=\; \tilde\Theta\!\Bigl(\frac{\sqrt r}{\varepsilon_2^2}
\;+\; \frac{r}{\log r}\cdot\max\Bigl\{\frac{\varepsilon_1}{\varepsilon_2^2},
\Bigl(\frac{\varepsilon_1}{\varepsilon_2^2}\Bigr)^2\Bigr\}\Bigr)\]
samples \autocite{canonne2022tolerance}. In particular, for constant
tolerances \(0<\varepsilon_1<\varepsilon_2\) the complexity is
\(\tilde\Theta(r/\log r)\) --- near-linear --- whereas the non-tolerant
case \(\varepsilon_1=0\) needs only \(\Theta(\sqrt r/\varepsilon_2^2)\)
\autocite{paninski2008coincidence,valiant2017automatic}. See also
\autocite{valiant2011unseen,jiao2018l1} for the \(\Omega(r/\log r)\)
lower bound at constant tolerance and the \(\ell_1\)-estimation
constants.
\end{quote}

Because §B.4 gives \(a>0\) (a \(\Theta(1)\) ball, not a point), our
testing problem has constant tolerances \(\varepsilon_1=a\),
\(\varepsilon_2=b\) and falls in the \emph{tolerant} regime: with
\(r=\Theta(n)\) types, any tester needs \(\tilde\Theta(n/\log n)\)
samples. The \(\Theta(\sqrt r)\) non-tolerant escape does not apply ---
and this is not an artifact but a design necessity, since demanding
consistency only at exactly-correct advice (\(a=0\)) would make the
consistency guarantee vacuous.

\section{Assembling the theorem}

Chaining §B.2--§B.5 on the family of §B.4 with constant
\(\theta,\varepsilon\) (hence \(\rho_{\mathrm{base}}=1-\Theta(1)\),
stakes \(\delta,\Delta=\Theta(1)\), tolerances \(a,b=\Theta(1)\),
support \(r=\Theta(n)\)):

\begin{enumerate}
\def\labelenumi{\arabic{enumi}.}
\tightlist
\item
  Suppose \(A\in\mathcal A_k\) with \(k=o(n/\log n)\) is
  \((1-o(1))\)-consistent and \((\rho_{\mathrm{base}}-o(1))\)-robust on
  the family.
\item
  By Lemma B.4, its decision solves the \((a,b)\)-tolerant
  identity-testing problem over support \(\Theta(n)\) with error
  \(o(1)\) from \(k\) samples.
\item
  By Theorem B.5, any such tester needs \(\tilde\Theta(n/\log n)\)
  samples --- more than \(k\) provides. The witness pair guaranteed by
  the lower bound consists of two type distributions \(p_G\) (at
  \(\ell_1(p_G,q)\le a\)) and \(p_{\mathrm{Bd}}\) (at
  \(\ell_1(p_{\mathrm{Bd}},q)\ge b\)) whose length-\(k\) prefix laws
  satisfy \(\gamma_k=o(1)\); by the affine law (Lemma B.3), \emph{any}
  pair at those distances realizes the gain \(\delta\) and loss
  \(\Delta\) respectively.
\item
  Lemma B.1 applied to \((p_G,p_{\mathrm{Bd}})\) with \(\gamma_k=o(1)\)
  forces \(\eta_c+\eta_r\ge 1-o(1)\) --- contradicting step 1.
\end{enumerate}

\begin{quote}
\textbf{Theorem B.6 (impossibility; the theorem of Chapter 9).} For
every test budget \(k=o(n/\log n)\) there is a known-i.i.d. matching
family with \(r=\Theta(n)\) types and baseline strength
\(\rho_{\mathrm{base}}=1-\Theta(1)\) on which no test-and-fallback
algorithm \(A\in\mathcal A_k\) --- deciding by any rule on its prefix
--- is simultaneously \((1-o(1))\)-consistent and
\((\rho_{\mathrm{base}}-o(1))\)-robust. The achievable (consistency,
robustness) region collapses to the baseline point
\((\rho_{\mathrm{base}},\rho_{\mathrm{base}})\).
\end{quote}

Combined with the easy strong-baseline side --- few types
\(\Rightarrow\) near-perfect matching \(\Rightarrow\) upside
\(\delta\to0\), nothing to capture --- this is the scissors of Figure
9.1: the upside exists only where the tolerant test is infeasible.

\section{Status, the remaining step, and a negative
finding}

\textbf{Proved here.} Lemma B.1 (complete proof, §B.2); Lemma B.2
(complete proof plus three-decimal numerical verification, §B.3); Lemma
B.3 (complete proof plus verification, §B.4); Lemma B.4 (proof sketch
whose completion is bookkeeping on the two conditioning displays of
Lemma B.1). Theorem B.5 is a cited, exact, modern result --- the
load-bearing external ingredient --- and §B.4 establishes that our
problem sits on its tolerant (theorem-preserving) side.

\textbf{The remaining step (stated plainly).} Step 3 of §B.6 requires
\emph{exhibiting} the witness pair \((p_G,p_{\mathrm{Bd}})\) from the
tolerant-testing lower-bound construction and checking it lands at
\(\ell_1\le a\) and \(\ell_1\ge b\) within the cell parameterization of
§B.4. Because the affine law makes any pair at those distances realize
the stakes, this is a routine verification rather than new mathematics;
it, plus a final end-to-end review, remains open at the time of writing
(see §9.5). Until it is closed, Theorem B.6 should be read with the
status stated in Chapter 9.

\textbf{A negative finding worth recording.} A \emph{low-dimensional}
version of the construction --- all cells sharing one global bias
parameter, with the good/bad scenarios differing by a single global
shift --- provably fails: a length-\(k\) prefix estimates one shared
parameter to \(\pm1/\sqrt{\theta k}\), so the scenarios become
distinguishable at any \(\delta=\omega(\sqrt{\theta/k})\) and only a
vacuous impossibility survives. The strength of the theorem genuinely
requires spreading the advice error across \(\Theta(n)\) independent
cells, each perturbed below the per-type sampling resolution --- the
dimension is essential, which is also why the theorem's strong form is
stated for the \(r=\Theta(n)\) regime and a constant-\(r\) version is
left open (§9.5).
